% Reusable undated syllabus, adapted from the supplied Google Doc.
% Source: https://docs.google.com/document/d/1rP35xUjv5yG_3O_kjvSFSWJbrn52myI_LdM4LQRM4iA/edit
% Original text and DOCX snapshots are in sources/.
% Bibliographic years and article titles are intentionally unchanged.
% Class numbering follows the source, including Class 3.2 (no Class 3.1).
\documentclass[11pt,letterpaper]{article}
\usepackage[T1]{fontenc}
\usepackage[utf8]{inputenc}
\usepackage{lmodern}
\usepackage[margin=0.8in,headheight=14pt,headsep=15pt,footskip=24pt]{geometry}
\usepackage{microtype}
\usepackage{xcolor}
\usepackage{enumitem}
\usepackage{titlesec}
\usepackage{needspace}
\usepackage{tabularx}
\usepackage{fancyhdr}
\usepackage{xurl}
\usepackage{hyperref}
\definecolor{maroon}{HTML}{800000}
\definecolor{linkblue}{HTML}{174E79}
\hypersetup{colorlinks=true,urlcolor=linkblue,linkcolor=maroon,
  pdftitle={Lab and Field Experiments in Comparative Politics and Policy},
  pdfauthor={Molly Offer-Westort},pdfsubject={PLSC 48110/28110 syllabus}}
\urlstyle{same}
\setlength{\parindent}{0pt}
\setlength{\parskip}{6pt}
\setlength{\emergencystretch}{2em}
\raggedbottom
\clubpenalty=10000
\widowpenalty=10000
\pagestyle{fancy}
\fancyhf{}
\fancyhead[L]{\small\textcolor{maroon}{PLSC 48110/28110}}
\fancyhead[R]{\small Lab and Field Experiments}
\fancyfoot[R]{\small\thepage}
\renewcommand{\headrulewidth}{0.3pt}
\titleformat{\section}{\large\bfseries\color{maroon}}{}{0pt}{}
\titlespacing*{\section}{0pt}{16pt}{6pt}
\titleformat{\subsection}{\normalsize\bfseries}{}{0pt}{}
\titlespacing*{\subsection}{0pt}{11pt}{4pt}
\setlist[itemize]{leftmargin=1.25em,itemsep=5pt,topsep=3pt,parsep=0pt}
\setlist[itemize,2]{leftmargin=1.4em,itemsep=3pt}
\setlist[itemize,3]{leftmargin=1.4em,itemsep=2pt}
\newcommand{\courseweek}[2]{\Needspace{9\baselineskip}\section*{Week #1: #2}}
\newcommand{\class}[2]{\Needspace{6\baselineskip}\subsection*{Class #1}\nopagebreak #2\par\nopagebreak\textbf{Readings}\nopagebreak}
\newcommand{\blankfield}{\rule{1.75in}{0.3pt}}
\newcommand{\showfield}[1]{\if\relax\detokenize{#1}\relax\blankfield\else#1\fi}
% Fill in these four fields for a particular offering. Empty values render as blank lines.
\newcommand{\CourseMeetingTime}{}
\newcommand{\CourseLocation}{}
\newcommand{\InstructorOfficeHours}{}
\newcommand{\InstructorOfficeLocation}{}
% The booking link is retained from the source; update it as needed.
\newcommand{\OfficeHoursBookingURL}{https://calendly.com/mollyow/office-hours}

\begin{document}
\thispagestyle{plain}
{\small\bfseries\color{maroon} UNIVERSITY OF CHICAGO\par}
\vspace{4pt}
{\LARGE\bfseries Lab and Field Experiments in\\[2pt]Comparative Politics and Policy\par}
\vspace{4pt}
{\large PLSC 48110/28110\par}
\vspace{5pt}
\textbf{Instructor:} Dr. Molly Offer-Westort, \href{mailto:mollyow@uchicago.edu}{mollyow@uchicago.edu}

\begin{tabularx}{\linewidth}{@{}lXlX@{}}
\textbf{Meetings:} & \expandafter\showfield\expandafter{\CourseMeetingTime} & \textbf{Location:} & \expandafter\showfield\expandafter{\CourseLocation}\\[5pt]
\textbf{Office hours:} & \expandafter\showfield\expandafter{\InstructorOfficeHours} & \textbf{Office:} & \expandafter\showfield\expandafter{\InstructorOfficeLocation}
\end{tabularx}

\textbf{Office-hours booking:} \href{\OfficeHoursBookingURL}{Book an appointment}

\section*{Course description and objectives}


Over the last several decades in what has become known as an ``experimental revolution,'' social scientists have increasingly used experimental designs to study the effects of interventions and policies. The growing popularity of experimental methodology gives researchers leverage on causal questions, and bringing these designs to the field helps researchers and policy makers to take their research designs closer to the real social phenomena they seek to understand. In this course, we will study how these methods have been used to study politics and policy, with examples from Africa, Latin America, and India. We will ask what types of questions experiments can answer, and what types of things we can manipulate and measure. We will review design considerations, such as methods to account for treatment spillover, and randomization procedures. We will also read critiques of experimental methods and field experiments; students will discuss ethical considerations, and will consider tradeoffs to limitations of experimental designs in field settings. While we will carry out some calculations, this course will not be oriented towards statistical programming; however, undergraduate class participants should have taken PLSC 26969 or obtain instructor consent.

\section*{Assignments}

Each student will be responsible for co-leading discussion on two of the second meetings of the week, beginning the third week of the course. For these meetings, discussion leaders will jointly provide a 20-30 minute total high-level overview of the papers assigned, and will develop a list of discussion questions based on the course's running substantive and methodological themes. Co-discussion leaders may divide papers among themselves or may collaborate; presentation time is combined total across all papers and discussion leaders. Discussion questions should be submitted the evening before class.

The final project will be a pre-analysis plan for a hypothetical lab or field experiment (not survey experiment) relevant to comparative politics. In the sixth week of the class, students should submit one page overviews of the proposed research design, and will meet with the instructor to discuss.

\textbf{Grade composition}

\begin{itemize}

\item Leading class discussions: 25\% for each of two classes

\item Final project: 40\%

\item Participation: 10\%

\end{itemize}

\section*{Accommodations}

Please reach out to me directly if you would like to request accommodations for the course to better facilitate your learning. Student Disability Services (\url{https://disabilities.uchicago.edu/}) is also available to provide you resources and support, and may provide approval for specific academic accommodations. If you or your household are affected by the ongoing pandemic in a way that affects your ability to participate in or attend class, please reach out to me as well. Informing me in a timely manner will help me to ensure accommodations are met and I am able to implement an appropriate assessment of your learning.

\section*{Academic honesty}

All students are expected to be familiar with and follow the University's academic honesty policies (\url{https://studentmanual.uchicago.edu/academic-policies/academic-honesty-plagiarism/}). For this course, you may use generative artificial intelligence tools such as ChatGPT for brainstorming, but you may not directly use any text generated from such tools in any assignments. Note that artificial intelligence may ``hallucinate'' information, so exercise extra caution and verify any information AI provides with reliable sources and proper citation practices
(see discussions on AI hallucination in: \url{https://www.nytimes.com/2023/03/29/technology/ai-chatbots-hallucinations.html}, \url{https://www.nytimes.com/2023/11/06/technology/chatbots-hallucination-rates.html}).

\section*{Course schedule}

\courseweek{1}{Course introduction}

\class{1.1}{Class overview. Social science questions we ask and answer with experiments. Randomization inference.}

\begin{itemize}

\item None.

\item For reference (you don't have to read these)

\begin{itemize}

\item Holland, Paul W. (1986). Statistics and causal inference. \emph{Journal of the American Statistical Association, 81}(396), 945-960. \url{https://doi.org/10.1080/01621459.1986.10478354}

\item King, Keohane, and Verba (1994). Designing Social Inquiry: \emph{Scientific Inference in Qualitative Research. }Princeton University Press. Chapter 1. \url{https://doi.org/10.1515/9781400821211-002}

\end{itemize}

\end{itemize}

\clearpage
\courseweek{2}{Introduction to estimation, and survey experiments}

\class{2.1}{Overview of the recent history of field experiments for policy and politics, and some prospective considerations.}

\begin{itemize}

\item Levitt, S. D., \& List, J. A. (2008). Field experiments in economics: The past, the present, and the future. \emph{European Economic Review}, 53(1), 1-18. \url{https://doi.org/10.1016/j.euroecorev.2008.12.001}

\item Deaton, A. (2020). Randomization in the tropics revisited: a theme and eleven variations (No. w27600). National Bureau of Economic Research.  \url{https://doi.org/10.3386/w27600}

\item Banerjee, A., Duflo, E., Glennerster, R., \& Kinnan, C. (2015). The miracle of microfinance? Evidence from a randomized evaluation. \emph{American economic journal: Applied economics}, 7(1), 22-53.  \url{http://dx.doi.org/10.1257/app.20130533}

\end{itemize}

\class{2.2}{Estimation and inference for randomized experiments.}

\begin{itemize}

\item Gerber, A. S., \& Green, D. P. (2012). \emph{Field experiments: Design, analysis, and interpretation}. WW Norton. Chapter 2: Causal Inference and Experimentation.

\item Athey, S., \& Imbens, G. W. (2017). The econometrics of randomized experiments. In \emph{Handbook of economic field experiments} (Vol. 1, pp. 73-140). North-Holland. \url{https://doi.org/10.1016/bs.hefe.2016.10.003}. Parts 1-4 only.

\end{itemize}

\courseweek{3}{Lab and lab-in-the-field experiments}

\class{3.2}{Applications of behavioral games and lab-in-the-field experiments.}

\begin{itemize}

\item Henrich, J., Boyd, R., Bowles, S., Camerer, C., Fehr, E., Gintis, H., \& McElreath, R. (2001). In search of homo economicus: behavioral experiments in 15 small-scale societies. \emph{American Economic Review}, 91(2), 73-78. \url{https://doi.org/10.1257/aer.91.2.73}

\item Habyarimana, J., Humphreys, M., Posner, D. N., \& Weinstein, J. M. (2007). Why does ethnic diversity undermine public goods provision?. \emph{American political science review}, 101(4), 709-725. \url{https://doi.org/10.1017/S0003055407070499}

\item Baldassarri, D., \& Grossman, G. (2011). Centralized sanctioning and legitimate authority promote cooperation in humans. \emph{Proceedings of the National Academy of Sciences}, 108(27), 11023-11027. \url{https://doi.org/10.1073/pnas.1105456108}

\end{itemize}

\clearpage
\courseweek{4}{Estimation considerations}

\class{4.1}{Experiments and covariate adjustment.}

\begin{itemize}

\item Athey, S., \& Imbens, G. W. (2017). The econometrics of randomized experiments. In \emph{Handbook of economic field experiments} (Vol. 1, pp. 73-140). North-Holland. \url{https://doi.org/10.1016/bs.hefe.2016.10.003}. Parts 5-6 only.

\item Lin, W. (2012). Regression Adjustment in Randomized Experiments: Is the Cure Really Worse Than the Disease?. Blog post.

\begin{itemize}

\item \href{https://blogs.worldbank.org/impactevaluations/regression-adjustment-in-randomized-experiments-is-the-cure-really-worse-than-the-disease}{Part I}

\item \href{https://blogs.worldbank.org/impactevaluations/guest-post-by-winston-lin-regression-adjustment-in-randomized-experiments-is-the-cure-really-worse-0}{Part II}

\end{itemize}

\end{itemize}

\class{4.2}{Measurement and results presentation.}

\begin{itemize}

\item Olken, B. A. (2010). Direct democracy and local public goods: Evidence from a field experiment in Indonesia. \emph{American political science review}, 104(2), 243-267. \url{https://doi.org/10.1017/S0003055410000079}

\item Wantchekon, L. (2003). Clientelism and voting behavior: Evidence from a field experiment in Benin. \emph{World politics}, 55(3), 399-422. \url{https://doi.org/10.1353/wp.2003.0018}

\item Anderson, S. E., Buntaine, M. T., Liu, M., \& Zhang, B. (2019). Non-governmental monitoring of local governments increases compliance with central mandates: a national-scale field experiment in China. \emph{American Journal of Political Science}, \emph{63}(3), 626-643. \url{https://doi.org/10.1111/ajps.12428}

\end{itemize}

\courseweek{5}{Design documents and pre-analysis plans}

\class{5.1}{Writing a design document or pre-analysis plan.}

\begin{itemize}

\item Gerber, A. S., \& Green, D. P. (2012). \emph{Field experiments: Design, analysis, and interpretation}. WW Norton. Chapter 13: Writing an Experimental Proposal, Research Report, and Journal Article.

\item United States. National Commission for the Protection of Human Subjects of Biomedical, \& Behavioral Research. (1978). \emph{The Belmont report: ethical principles and guidelines for the protection of human subjects of research}. Department of Health, Education, and Welfare, National Commission for the Protection of Human Subjects of Biomedical and Behavioral Research. \url{https://www.hhs.gov/ohrp/regulations-and-policy/belmont-report/read-the-belmont-report/index.html}

\Needspace{10\baselineskip}
\item Design documents (for reference, you can skim these):

\begin{itemize}

\item Young, L. E. (2015). The psychology of political risk: Replication and extension. \url{https://doi.org/10.17605/OSF.IO/WGJZ3}

\item Gaikwad, N., \& Nellis, G. (2019). Field Experiment on the Political Empowerment of Internal Migrants in India. \url{https://doi.org/10.17605/OSF.IO/7VTQH} Also read supplementary information sections K and L, available at \url{https://doi.org/10.7910/DVN/G1JCKK}

\end{itemize}

\end{itemize}

\class{5.2}{Applications from design documents, ethical considerations with human subjects research.}

\begin{itemize}

\item Young, L. E. (2019). The psychology of state repression: Fear and dissent decisions in Zimbabwe. \emph{American Political Science Review}, 113(1), 140-155. \url{https://doi.org/10.1017/S000305541800076X}

\item Gaikwad, N., \& Nellis, G. (2021). Overcoming the political exclusion of migrants: Theory and experimental evidence from India. \emph{American Political Science Review}, 115(4), 1129-1146. \url{https://doi.org/10.1017/S0003055421000435}

\end{itemize}

\courseweek{6}{Design considerations: power and randomization}

\class{6.1}{Power calculations and randomization procedures.}

\begin{itemize}

\item Athey, S., \& Imbens, G. W. (2017). The econometrics of randomized experiments. In \emph{Handbook of economic field experiments} (Vol. 1, pp. 73-140). North-Holland. \url{https://doi.org/10.1016/bs.hefe.2016.10.003}. Part 7 only.

\item JPAL \emph{Power calculations}.  \url{https://www.povertyactionlab.org/resource/power-calculations}

\end{itemize}

\class{6.2}{Different randomization procedures.}

\begin{itemize}

\item Blair, R. A., Moscoso-Rojas, M., Castillo, A. V., \& Weintraub, M. (2022). Preventing Rebel Resurgence after Civil War: A Field Experiment in Security and Justice Provision in Rural Colombia. \emph{American Political Science Review}, 116(4), 1258-1277. \url{https://doi.org/10.1017/S0003055422000284}

\item Blattman, C., Duncan, G., Lessing, B., \& Tobon, S. (2022). State-building on the margin: An urban experiment in Medellín (No. w29692). \emph{National Bureau of Economic Research}. \url{https://doi.org/10.31235/osf.io/3bncz}

\end{itemize}

\Needspace{3\baselineskip}\textbf{Assignment: Write up one-page research proposal. Meet with instructor this week and next to discuss.}

\courseweek{7}{Spillovers}

\class{7.1}{Estimation and design accounting for spillovers.}

\begin{itemize}

\item Gerber, A. S., \& Green, D. P. (2012). \emph{Field experiments: Design, analysis, and interpretation}. WW Norton. Chapter 8: Interference between Experimental Units.

\end{itemize}

\class{7.2}{Designs to detect spillovers.}

\begin{itemize}

\item Miguel, E., \& Kremer, M. (2004). Worms: identifying impacts on education and health in the presence of treatment externalities. \emph{Econometrica}, 72(1), 159-217. \url{https://doi.org/10.1111/j.1468-0262.2004.00481.x}

\begin{itemize}

\item Additional unrequired content on the ``Worm Wars'' regarding replication issues with Miguel \& Kremer:

\begin{itemize}

\item \href{https://www.vox.com/2015/7/24/9031909/worm-wars-explained}{Vox article}

\item \href{https://open.spotify.com/episode/7ujpdCPK5hKPgUYDGV87aV?si=nmXeJSqnR0aRlF5S3gMP9g}{Maintenance Phase podcast episode}

\end{itemize}

\end{itemize}

\item Wilke, A. M., Green, D. P., \& Cooper, J. (2020). A placebo design to detect spillovers from an education--entertainment experiment in Uganda. \emph{Journal of the Royal Statistical Society: Series A (Statistics in Society)}, 183(3), 1075-1096. \url{https://doi.org/10.1111/rssa.12571}

\end{itemize}

\courseweek{8}{Instrumental variables}

\class{8.1}{Instrumental variables estimation.}

\begin{itemize}

\item Angrist, J. D., \& Krueger, A. B. (2001). Instrumental variables and the search for identification: From supply and demand to natural experiments. \emph{Journal of Economic perspectives}, 15(4), 69-85. \url{https://doi.org/10.1257/jep.15.4.69}

\item Athey, S., \& Imbens, G. W. (2017). The econometrics of randomized experiments. In \emph{Handbook of economic field experiments} (Vol. 1, pp. 73-140). North-Holland. \url{https://doi.org/10.1016/bs.hefe.2016.10.003}. Part 9.

\end{itemize}

\clearpage
\class{8.2}{Applications with IV.}

\begin{itemize}

\item De La O, A. L. (2013). Do conditional cash transfers affect electoral behavior? Evidence from a randomized experiment in Mexico. \emph{American Journal of Political Science}, 57(1), 1-14.  \url{https://doi.org/10.1111/j.1540-5907.2012.00617.x}

\item Cameron, L., Erkal, N., Gangadharan, L., \& Meng, X. (2013). Little emperors: behavioral impacts of China's One-Child Policy. \emph{Science}, 339(6122), 953-957. \url{https://doi.org/10.1126/science.1230221}

\end{itemize}

\courseweek{9}{Heterogeneity and external validity}

\class{9.1}{Heterogeneous treatment effects and alternative experimental designs.}

\begin{itemize}

\item Athey, S., \& Imbens, G. W. (2017). The econometrics of randomized experiments. In \emph{Handbook of economic field experiments} (Vol. 1, pp. 73-140). North-Holland. \url{https://doi.org/10.1016/bs.hefe.2016.10.003}. Parts 10-11.

\end{itemize}

\class{9.2}{External validity.}

\begin{itemize}

\item Banerjee, A., Duflo, E., Goldberg, N., Karlan, D., Osei, R., Parienté, W., ... \& Udry, C. (2015). A multifaceted program causes lasting progress for the very poor: Evidence from six countries. \emph{Science}, 348(6236), 1260799. \url{https://doi.org/10.1126/science.1260799}

\item Abaluck, J., Kwong, L. H., Styczynski, A., Haque, A., Kabir, M. A., Bates-Jefferys, E., ... \& Mobarak, A. M. (2022). Impact of community masking on COVID-19: A cluster-randomized trial in Bangladesh. \emph{Science}, \emph{375}(6577), eabi9069. \url{https://doi.org/10.1126/science.abi9069}

\item For reference: blog posts on meta-analyses

\begin{itemize}

\item Simmons, J. P., Nelson, L. D., \&  Simonsohn, U. (2019). Meaningless Means Series. \emph{Data Colada}. Blog posts. \url{https://datacolada.org/meaningless}

\end{itemize}

\item Not required: Meager, R. (2022). "Aggregating distributional treatment effects: A Bayesian hierarchical analysis of the microcredit literature." \emph{American Economic Review} 112.6: 1818-1847. \url{https://doi.org/10.1257/aer.20181811}

\end{itemize}

\Needspace{3\baselineskip}\textbf{Pre-analysis plans due on the last day of class.}

\clearpage
\section*{Other readings}

\begin{itemize}

\item Aronow, P. M., \& Samii, C. (2017). Estimating average causal effects under general interference, with application to a social network experiment. \emph{The Annals of Applied Statistics}, 11(4), 1912-1947. \url{https://doi.org/10.1214/16-AOAS1005}

\item Athey, S., Bergstrom, K., Hadad, V., Jamison, J. C., Ozler, B., Parisotto, L., \& Sama, J. D. (2021). Shared Decision-Making: Can Improved Counseling Increase Willingness to Pay for Modern Contraceptives?. Research Papers 3987, Stanford University, Graduate School of Business. \url{https://doi.org/10.1596/1813-9450-9777}

\item Baird, S., Bohren, J. A., McIntosh, C., \& Özler, B. (2018). Optimal design of experiments in the presence of interference. \emph{Review of Economics and Statistics}, 100(5), 844-860. \url{https://doi.org/10.1162/rest_a_00716}

\item Bowles, S., \& Gintis, H. (2004). Homo Economicus and Zoon Politikon.

\item Cameron, L., Meng, X., \& Zhang, D. (2019). China's sex ratio and crime: Behavioural change or financial necessity?. \emph{The Economic Journal}, 129(618), 790-820. \url{https://doi.org/10.1111/ecoj.12561}

\item De Figueiredo, M. F., Hidalgo, F. D., \& Kasahara, Y. (2011). When do voters punish corrupt politicians? Experimental evidence from Brazil. \emph{Unpublished manuscript}, UC Berkeley.

\item Fehr, E., \& Fischbacher, U. (2002). Why social preferences matter--the impact of non-selfish motives on competition, cooperation and incentives. \emph{The economic journal}, 112(478), C1-C33.  \url{https://doi.org/10.1111/1468-0297.00027}

\item Gonzalez-Ocantos, E., De Jonge, C. K., Meléndez, C., Osorio, J., \& Nickerson, D. W. (2012). Vote buying and social desirability bias: Experimental evidence from Nicaragua. American \emph{Journal of Political Science}, 56(1), 202-217.

\item Ichino, N., \& Schündeln, M. (2012). Deterring or displacing electoral irregularities? Spillover effects of observers in a randomized field experiment in Ghana. \emph{The Journal of Politics}, 74(1), 292-307. \url{https://doi.org/10.1017/S0022381611001368}

\item Mansuri, G., \& Rao, V. (2013). Can participation be induced? Some evidence from developing countries. \emph{Critical Review of International Social and Political Philosophy,} 16(2), 284-304. \url{https://doi.org/10.1080/13698230.2012.757918}

\item Reinikka, R., \& Svensson, J. (2005). Fighting corruption to improve schooling: Evidence from a newspaper campaign in Uganda. \emph{Journal of the European economic association}, 3(2-3), 259-267. \url{https://www.jstor.org/stable/40004969}

\end{itemize}

\end{document}
